\section{The comparison error on the unhit region}
\label{sec:seed-unhit-comparison}
We compare the actual cut kernels with the full linear kernels at the
identity. The linear cancellation uses the derivative of the actual
rooted predictor. Its gradient correction disappears only after the
coarse curl is applied. Throughout this section $L$ and the patch
are fixed; constants may depend on them but are uniform in the period,
the truncation length, and the selected tags.

\subsection{The full linear identities}
Let $B_f,B_c$ be the two stacked curvature maps: their first components
are $\sqrt{a_f}$ and $\sqrt{a_c}$ times the identity, and their
second components are the full path-root kernels. Path coefficients
are aligned with these root kernels by summing all labels with the
same displacement and target orientation. Absolute convergence in
\eqref{eq:guide-path-input} permits periodization and grouping.
Consequently the linearizations of the two plaquette fields are
\[
 X_f^0=B_f d_f a,\qquad X_c^0=B_c d_c b,
\]
where $a,b$ are fine and coarse logarithmic bond fields and $d_f,d_c$
are the discrete exterior derivatives from bonds to plaquettes.
Let $F_0,I_0,N_0$ be the full real kernels corresponding to
$\Phi,\Phi',N$ at the identity. The flat stack identities needed
here are
\begin{equation}\label{eq:unhit-flat-stack}
 \begin{split}
 F_0^*F_0+I_0^*I_0+N_0^*N_0&=1,\\
 \sqrt p\,I_0^*&=B_cd_c,\\
 F_0^*B_fd_f&=\sqrt p\,I_0^*\mathcal A.
 \end{split}
\end{equation}
Here $\mathcal A$ is the unrooted linear averaging map. These are
identities of full finite matrices in physical counting coordinates,
for the same step and the same two stages. We state their exact role
because none of them is asserted for the transported cut kernels.
Their production, together with the aligned path expansion, is the
flat-stack input to this comparison.

The derivative $\mathcal Q$ of the actual predictor has the form
$\mathcal Q=\mathcal A+d_c^{(0)}\mathcal R$, where
$d_c^{(0)}$ sends site fields to bond gradients. Thus, since
$d_cd_c^{(0)}=0$,
\[
 \sqrt p\,I_0^*\mathcal Q=B_cd_c\mathcal A.
\]
The discrete identity $d_cd_c^{(0)}=0$ follows by cancellation of
the four endpoint differences around each plaquette. With
$Z^0=\sqrt p\,(b-\mathcal Qa)$, the last two identities in
\eqref{eq:unhit-flat-stack} therefore give
\begin{equation}\label{eq:unhit-linear-zero}
 F_0^*X_f^0+I_0^*Z^0-X_c^0=0.
\end{equation}
The first identity in \eqref{eq:unhit-flat-stack} also gives
$(1-F_0^*F_0-I_0^*I_0-N_0^*N_0)w=0$ for every doubled coarse
field $w$. No closedness assumption on $w$ or removal of harmonic
modes is used.

\subsection{Uniform Taylor bounds for local reads}
For a real logarithmic field $a$, put $U_e=\exp(a_e)$ and use its
supremum norm. On a fixed small group chart,
\[
 r(\exp v)\le C\|v\|,\qquad
 \|\log g\|\le C r(g),\qquad \exp(\log g)=g.
\]
These follow by integrating the first derivative of the analytic
exponential and local logarithm on smaller closed balls. Likewise,
Taylor's formula for the product of the four exponentials around a
plaquette gives
\begin{equation}\label{eq:unhit-face-taylor}
 \|F(U_p)-(d a)_p\|\le C\|a\|_\infty^2,
 \qquad \|(d a)_p\|\le C\|a\|_\infty.
\end{equation}
Indeed differentiating the product at zero gives the signed sum of
the four logarithms; the derivative of the curvature coordinate is
the identity on the Lie algebra. Bound the second derivative on the
fixed four-slot ball and integrate the remainder. Repeated periodic
bonds simply repeat coordinates in that four-slot map and do not
change the constant.

For the actual predictor we use its local analytic law at the
identity, for each allowed tag, with derivative $\mathcal Q$ above.
The law is part of the branch construction: it includes uniqueness
of the normalized minimizing frame on the branch where that frame
is used. It implies the uniform estimate
\begin{equation}\label{eq:unhit-predictor-taylor}
 \|Z-Z^0\|_\infty\le C h^2,\qquad
 \|Z^0\|_\infty\le Ch,
 \quad h=\max\{\|a\|_\infty,\|b\|_\infty\},
\end{equation}
with a radius independent of the torus. Here is why finite-dimensional
Taylor estimates give that uniformity. A predictor bond reads only
two frame patches, their coordinate paths, and their length-$L$
connectors. This is a fixed finite set of cover bond anchors. Choose
one reference torus large enough that these anchors have their actual
cover lifts and its two coarse endpoints are distinct. Transplant
the local fine logarithms to it by reading the original periodic
field at each corresponding cover anchor. Transplant the coarse
logarithm at the tested bond as a constant field, and copy the actual
two frame tags and connector tag. The transplanted sup norm does not
increase. All path products, frame scores, selected frames and
connectors agree. On a minimizing branch, equality of scores gives
equality of the selected normalized frame by its uniqueness; on the
other branch the chosen axial paths agree directly.

There are only finitely many tags and tested directions on this
reference torus. Take the minimum of their positive Taylor radii and
the maximum of their Taylor bounds. Equality of the transplanted
reads along every ray $v\mapsto v(a,b)$ gives equality of their
derivatives at zero as well, so these bounds prove
\eqref{eq:unhit-predictor-taylor} on every torus, including short
ones. This proves the uniformization; the stated analytic branch law
and identification of its derivative are required inputs.

\subsection{Truncation and transport errors}
For cutoff $R$ in fine units define
\[
 \omega(R,h)=(1+R)h+e^{-\tau R},\qquad \tau>0.
\]
The kernel strip width here is strictly positive, $\eta>0$.
The value of $\tau$ can be decreased a finite number of times below.
If all logarithms have norm at most $h$ in the chart, every link has
chord at most $Ch$, and a path of length $k$ has chord at most $Ckh$
by the product triangle inequality. On the colour space,
\begin{equation}\label{eq:unhit-adjoint}
 \|\operatorname{Ad}(g)-1\|\le2D r(g).
\end{equation}
Here $D$ is the fixed colour normalization. To see this, write
$gXg^*-X=(g-1)Xg^*+X(g^*-1)$, use the Frobenius product bound and
unitary invariance, then divide by the colour normalization.

Consider a transported kernel term with cover displacement $z$ and
physical path length $k\le D_L(1+|z|_1)$, with $|z|_1$ the $\ell^1$
norm as in Section~\ref{sec:seed-guides}. This bound holds for the
actual paths: a narrow path has fine length $L|z|_1$, while the wide
path from residue $r$ to $Lz$ has length at most $L|z|_1+4L$.
If $k\le R$, \eqref{eq:unhit-adjoint} bounds the difference from
the scalar flat term by $C Rh|K_{ab}(z)|$. If $k>R$, the flat term
is omitted, and
\[
 1\le e^\eta e^{\eta|z|_1}e^{-\eta R/D_L}.
\]
Paying both cases with the weighted strip sum and using the same
row and column counting as \eqref{eq:guide-periodic-rows} proves
\begin{equation}\label{eq:unhit-kernel-near}
 \|F\|_{\rm row},\|F\|_{\rm col},
 \|F_0\|_{\rm row},\|F_0\|_{\rm col}\le C,
 \qquad
 \|F-F_0\|_{\rm row},\|F-F_0\|_{\rm col}
       \le C\omega(R,h).
\end{equation}
This holds for each of the three kernels, after increasing $C$ and
choosing $\tau\le\eta/D_L$. The type counts, including residues,
are absorbed after $L$ is fixed. Every cover term is estimated before
periodization, so winding paths cause no exception.

The same argument applies to the two plaquette fields. For a retained
path term, first use \eqref{eq:unhit-face-taylor}, then
\eqref{eq:unhit-adjoint} on its transport. Its error is at most
$C(1+R)h^2$ times its coefficient. For an omitted term use its full
linear value, bounded by $Ch$, and
$1\le e^{\sigma\ell}e^{-\sigma R/\kappa}$ when $\kappa\ell>R$, where $\kappa\in\{1,L\}$
is the stage factor of Section~\ref{sec:seed-guides}.
The weighted absolute path sum therefore gives
\begin{equation}\label{eq:unhit-field-near}
 \|X^0\|_\infty\le Ch,\qquad
 \|X-X^0\|_\infty\le Ch\omega(R,h)
\end{equation}
for both stages, with $\tau\le\min\{\sigma_f,\sigma_c/L\}$.
Absolute convergence justifies subtracting and summing the path
series, even if the set of labels below the length cutoff is infinite.
Equation \eqref{eq:unhit-predictor-taylor} has the same form since
$h\le\omega(R,h)$.

If $F,F_0$ satisfy \eqref{eq:unhit-kernel-near} and
$z,z^0$ satisfy \eqref{eq:unhit-field-near}, then
\[
 F^*z-F_0^*z^0=F^*(z-z^0)+(F-F_0)^*z^0
\]
has sup norm at most $2C^2h\omega(R,h)$. For an arbitrary bounded
field $w$, the same decomposition gives
\[
 \|(F^*F-F_0^*F_0)w\|_\infty
       \le2C^2\|w\|_\infty\omega(R,h).
\]
Insert these bounds in the two full linear cancellations. After
increasing $C$, the actual defect vectors satisfy
\begin{equation}\label{eq:unhit-global}
 \|\Phi^*X_f+(\Phi')^*Z-X_c\|_\infty\le Ch\omega(R,h),
 \qquad
 \|(1-\Phi^*\Phi-(\Phi')^*\Phi'-N^*N)w\|_\infty
       \le C\|w\|_\infty\omega(R,h).
\end{equation}
The argument keeps $w$ arbitrary; it requires no regularity of the
clipping indicators that will later define $w$.

\subsection{Restriction to a neighbourhood and summation}
A cut kernel entry has endpoints within distance $R$ in coarse
anchor units and reads only the transport path of that length.
The cut path field reads through radius $R+4$. A predictor bond
reads through a fixed patch radius $R_{\rm pred}$. Thus the first
defect at anchor $o$ depends only on fine and coarse bonds through
radius $3R+4+R_{\rm pred}$, using parent distance for fine bonds.
For example, the outer transpose moves the input anchor by at most
$R$, its kernel path adds at most $R$, and the input field adds
$R+4$ or the predictor radius. The second defect, as an operator
applied to a fixed $w$, reads its kernels through radius $2R$.
Equality of these reads gives equality term by term in each cut
series and finite matrix sum.

Suppose a coarse gauge $z$ makes the fine and coarse bond chords
at most $h$ on this neighbourhood. Take the local logarithms there
and set them to zero outside. Their global sup norms are at most
$Ah$ and their exponentials agree with the gauged fields on all
required reads. Apply \eqref{eq:unhit-global} to these extensions.
Replacing $h$ by $Ah$ changes only the constant, since for $A\ge1$,
$\omega(R,Ah)\le A\omega(R,h)$.

To remove the gauge, every transported plaquette term transforms
by $\operatorname{Ad}(z_o)$ at its anchor; the endpoint gauge
factors cancel. The same endpoint calculation for kernels gives
\[
 F'_{ij}=\operatorname{Ad}(z_i)F_{ij}\operatorname{Ad}(z_j)^{-1},
\]
with the parent anchor on fine rows. Taking physical adjoints
preserves this identity, because adjoint transport is orthogonal.
The actual predictor is covariant with the same tags under the
block-constant fine gauge. Hence both defect vectors transform by
$\operatorname{Ad}(z_o)$, provided the free field $w$ is transformed
at its anchors. Its norm is unchanged. Thus the local version of
\eqref{eq:unhit-global} holds for the original fields and for any
bounded $w$. In particular we may use the original clipped field
$v$ without changing its indicators in the extension argument.

Take $R=\lfloor B/100000\rfloor$ and
$h=C_LB^{10}T$ from \eqref{eq:tree-common-gauge}. For large $B$,
$3R+4+R_{\rm pred}\le\lfloor2B/5\rfloor$, so this estimate applies
to every unhit outer plaquette. As $\|v\|_\infty\le dt$, its two
comparison rows are bounded below by minus
\begin{equation}\label{eq:unhit-row-loss}
 \mathcal L=2dt\,Ch\omega(R,h)+dt\,Cdt\,\omega(R,h).
\end{equation}
There are two components in the linear row, while the factor $1/2$
in the quadratic row cancels its two components.

For completeness, the required delay is stronger than $h\to0$.
Write $H=x+c_L$, so $T=H^{-49/100}$ and
$t=P(H)H^{-1/2}\ge H^{-1/2}$. Any fixed power of
$B=\lceil(\log x)^2\rceil$ is smaller than any positive power of
$H$ eventually. Consequently, for every fixed integer $n$,
\[
 B^n h\longrightarrow0,\qquad
 B^n\frac{h^2}{t}\longrightarrow0,
\]
the latter because $T^2/t\le H^{-48/100}$. Also
$e^{-\tau R}\le H^{-1}$ eventually: $R$ grows quadratically in
$\log x$, whereas $\log H$ grows linearly. It follows that
\[
 B^n\omega(R,h)\longrightarrow0,\qquad
 B^n\frac{h\omega(R,h)}t\longrightarrow0.
\]
Expanding $\omega$ proves the second assertion from the $h^2/t$
bound and $e^{-\tau R}/t\le H^{-1/2}$. The outer region has at most
$NB^4n$ plaquettes for $n$ seeds. Given $\epsilon>0$, choose the
coupling threshold so that $NB^4\mathcal L\le\epsilon t^2$.
Summing \eqref{eq:unhit-row-loss} gives
\begin{equation}\label{eq:unhit-final}
 \sum_{p\text{ outer, unhit}}(L_p+D_p)\ge-\epsilon t^2 n.
\end{equation}
This is the unhit-region estimate used in the energy allocation.
Its remaining construction inputs are the full flat stack and aligned
path expansions, and the analytic law and derivative of the actual
predictor. The estimates, localization, common-gauge use and delay
from those inputs have all been given above.
